% ============================================================
\section{Cross-Block Commutator Defect (Observable/Shadow Mismatch)}
\label{sec:holographic-defect}
% ============================================================

\begin{quote}\small
\textbf{What this automates.}
When two operators do not commute, the ordering matters.
Under a horizon projection, the \emph{observed} commutator may differ from the commutator of \emph{observed} operators.
This section provides the language for ``holographic'' or ``bulk/boundary'' defects: mismatches between structure computed in the full space versus structure computed in the observable subspace.
\end{quote}

\subsection{Abstract pairing and ordering operators}

\begin{definition}[Pairing and ordering operators]
\label{def:pairing-ordering}\lean{PairingOrdering}
Let $\cH$ be a Hilbert space.
A \emph{pairing operator} $P:\cH \to \cH$ and an \emph{ordering operator} $O:\cH \to \cH$ are linear operators encoding two distinct structural operations on $\cH$.

No specific combinatorial or algebraic form is assumed; the operators are abstract.
\end{definition}

\paragraph{Motivation.}
In applications, $P$ might encode ``matching'' or ``coupling'' structure, while
$O$ encodes ``sequencing'' or ``precedence'' structure. The commutator $[P, O]$
measures the extent to which these structures interfere.

\subsection{Commutator defect under projection}

\begin{definition}[Commutator]
\label{def:commutator}\lean{commutator}
For operators $P, O:\cH \to \cH$, the \emph{commutator} is
\[
[P, O] := PO - OP.
\]
\end{definition}

\begin{definition}[Observed commutator]
\label{def:observed-commutator}\lean{observedCommutator}
The \emph{observed commutator} at horizon $h$ is
\[
[P, O]_h := \CCPi{h} [P, O] \CCPi{h} = \CCPi{h} (PO - OP) \CCPi{h}.
\]
\end{definition}

\begin{definition}[Commutator coherence defect]
\label{def:commutator-defect}\lean{commutatorCoherenceDefect}
The \emph{commutator coherence defect} is the coherence defect (Definition~\ref{def:coherence-defect}) applied to the commutator:
\[
\CCTau{h}{[P,O]}{x} := \CCPi{h} [P,O](x) - [P,O](\CCPi{h} x).
\]
This measures the mismatch between:
\begin{itemize}[itemsep=0.2em]
\item projecting the commutator's action on $x$, versus
\item computing the commutator's action on the projected state $\CCPi{h} x$.
\end{itemize}
\end{definition}

\paragraph{Two sources of mismatch.}
For the observed commutator $[P,O]_h$, there are two potential sources of
defect:
\begin{enumerate}[label=(\roman*)]
\item The commutator $[P,O]$ itself may be nonzero.
\item The coherence defect $\CCTau{h}{[P,O]}{x}$ may be nonzero even when
      $[P,O] \ne 0$.
\end{enumerate}
Full ``coherent commutativity'' requires both $[P,O] = 0$ and
$\CCTau{h}{[P,O]}{x} = 0$.

\subsection{Coherent stability}

\begin{definition}[Coherently stable state]
\label{def:coherently-stable}\lean{isCoherentlyStable}
A state $x \in \cH$ is \emph{coherently stable} (with respect to $P$, $O$, and $\CCPi{h}$) if
\[
[P, O]_h \, \CCPi{h} x = 0,
\]
i.e., the observed commutator annihilates the observable part of $x$.
\end{definition}

\paragraph{Interpretation.}
Coherent stability means the ordering ambiguity between $P$ and $O$ is
invisible to an observer at horizon $h$ when acting on $x$. The bulk structure
(full commutator) may be nontrivial, but its observable projection on $x$
vanishes.

\begin{proposition}[Coherent stability and defect vanishing]
\label{prop:coherent-stability}\lean{coherent_stability_observable}
If $x \in \cH_{\mathrm{obs}} = \CCPi{h} \cH$ (i.e., $\CCPi{h} x = x$), then $x$ is coherently stable if and only if
\[
\CCPi{h} [P, O] x = 0.
\]
\end{proposition}

\begin{proof}
For $x \in \cH_{\mathrm{obs}}$, we have $\CCPi{h} x = x$, so $[P,O]_h \CCPi{h} x = \CCPi{h} [P,O] \CCPi{h} x = \CCPi{h} [P,O] x$.
\end{proof}

\paragraph{No dimension selection claimed.}
This section provides only definitions and elementary consequences. Any claim
that coherent stability selects specific dimensions, subspaces, or critical
points would require additional structure and proof; such claims are not made
here.

\subsection{Cross-block decomposition}

\begin{proposition}[Cross-block commutator decomposition]
\label{prop:holographic-decomposition}\lean{crossBlockCommutatorDefect}
The commutator $[P,O]$ decomposes into observable and shadow contributions:
\[
[P,O] = [P,O]_h + \CCPibar{h} [P,O] \CCPi{h} + \CCPi{h} [P,O] \CCPibar{h} + \CCPibar{h} [P,O] \CCPibar{h}.
\]
The \emph{cross-block defect} is the sum of off-diagonal terms:
\[
\Xi_h^{\mathrm{cross}}[P,O] := \CCPibar{h} [P,O] \CCPi{h} + \CCPi{h} [P,O] \CCPibar{h}.
\]
(We write $\Xi$ rather than $\Delta$ to avoid collision with the increment projector $\CCDelta{h}$.)
\end{proposition}

\begin{proof}
Expand using $I = \CCPi{h} + \CCPibar{h}$ on both sides of $[P,O]$.
\end{proof}

\begin{quote}\small
\textbf{Intuition.}
The cross-block defect $\Xi_h^{\mathrm{cross}}$ captures the ``boundary'' terms: commutator contributions that cross between observable and shadow subspaces.
These are the terms lost when restricting attention to the observed commutator $[P,O]_h$ alone.
\end{quote}
